\documentclass[10pt]{article}

% graphicx package, useful for including eps, pdf, and raster graphics
\usepackage{graphicx}
\DeclareGraphicsExtensions{.png,.png,.jpg,.pdf}

% basic packages
\usepackage{fancyhdr}
\usepackage{color}
\usepackage{parskip}
\usepackage{float}
\usepackage{microtype}
\usepackage{url}
\usepackage{hyperref}
\usepackage{booktabs}
\usepackage{makecell}
\usepackage{multirow}
\usepackage{pbox}
\usepackage{enumitem}

% page layout
\usepackage[top=1in, bottom=1in, left=0.75in, right=0.75in, includefoot, heightrounded]{geometry}

% Useful for seeing the layout as you tweak it.
%\usepackage{showframe}

% fonts
\usepackage{fontspec}

\usepackage{polyglossia}
\setmainlanguage{vietnamese}
\setotherlanguage{english}

\setmainfont{DroidSans}[
  Extension   = .ttf,
  UprightFont = *,
  BoldFont    = *-Bold ]
\setmonofont{DroidSansMono}[
  Extension   = .ttf,
  UprightFont = * ]

\hypersetup{colorlinks,urlcolor=blue}
\urlstyle{same}

\setlength{\headheight}{1.1in}
\pagestyle{fancy}
\fancyhf{}

\lhead{\includegraphics[height=1in,clip]{logo}}
\rhead{\pbox[b][1in][c]{\textwidth}{
    \small
    \textbf{Tên người thực hiện xét nghiệm:}\\
    Northwest Genomics Center\\
    University of Washington School of Medicine\\
    Genome Sciences Box 355065\\
    3720 15th Ave NE\\
    Seattle, WA 98195-5065}}

\newcommand{\PageLine}{\rule{\textwidth}{0.25mm}}

\newcommand{\link}[2]{\href{#1}{#2}\footnote{#1}}

\begin{document}

\begin{center}
\Large
\textbf{Báo Cáo Xét Nghiệm SARS-CoV-2 (COVID-19) Bằng Phản Ứng Chuỗi Polymerase
(Polymerase Chain Reaction, PCR) Định Tính Của NWGC}
\end{center}

\bigskip

\begin{description}[font=\normalfont,align=left,labelwidth=14em]
\item [Tên Người Tham Gia] \textbf{\VAR{pat_name|e}}
\item [Ngày sinh] \textbf{\VAR{birth_date|e}}
\item [Mã định danh mẫu bệnh phẩm] \textbf{\VAR{qrcode|e}}
\item [Ngày Gửi Mẫu] \textbf{\VAR{collect_ts|e}}
\item [Ngày cung cấp kết quả]
  %- if status_code == "never-tested"
  \textbf{không bao giờ xét nghiệm}
  %- else
  \textbf{\VAR{result_ts|e}}
  %- endif
\item [Loại Mẫu Bệnh Phẩm] Cuốn mũi giữa, tự lấy mẫu
\end{description}

\PageLine

Kết Quả Xét Nghiệm SARS-CoV-2 (COVID-19) Bằng PCR Định Tính

%- if status_code == "negative"
\textbf{Không phát hiện}\\
%- elif status_code == "positive"
\textbf{Phát Hiện SARS-CoV-2  (dương tính)}\\
%- elif status_code == "inconclusive"
\textbf{Chưa có kết luận}\\
%- elif status_code == "never-tested"
\textbf{Không xét nghiệm được}\\
%- else
\VAR{0/0} %# poor man's raise
%- endif

\PageLine

Xin Lưu Ý:

\begin{itemize}
\item

  \textbf{Loại mẫu và xét nghiệm sử dụng trong dự án này được thực hiện như là
  một phần của chương trình giám sát y tế công cộng, với sự hợp tác cùng Quận
  King và các cơ quan khác.} Xét nghiệm mà quý vị đã nhận được không thể thay
  thế cho việc nhận chăm sóc y tế cá nhân từ nhà cung cấp dịch vụ chăm sóc sức
  khỏe.

\item

  \textbf{Chuỗi Kiểm Soát Mẫu Bệnh Phẩm:} mẫu bệnh phẩm xét nghiệm được tự lấy,
  tự dán nhãn và gửi đến Mạng Lưới Đánh Giá Về Vi-rút Corona của Seattle
  (Seattle Coronavirus Assessment Network - SCAN) qua đường bưu điện hoặc dịch
  vụ giao hàng. Mẫu bệnh phẩm không phải do chuyên gia y tế được cấp phép tiến
  hành thu thập, dán nhãn và vận chuyển. Điều này có thể ảnh hưởng đến kết quả
  xét nghiệm và không loại trừ khả năng là kết quả âm tính giả. Có thể cho phép
  xét nghiệm lâm sàng theo dõi tiếp.

\item

  \textbf{Thiết bị lấy mẫu bệnh phẩm:} que phết dịch tiết ở cuốn mũi giữa không
  phải là loại mẫu bệnh phẩm được phê duyệt theo Giấy Phép Sử Dụng Trong Trường
  Hợp Khẩn Cấp (Emergency Use Authorization - EUA) của Cơ Quan Quản Lý Dược Phẩm
  (Food and Drug Administration - FDA)/Trung Tâm Kiểm Soát và Phòng Ngừa Dịch
  Bệnh (Centers for Disease Control and Prevention - CDC). Do đó, kết quả không
  phát hiện/âm tính có thể sẽ không loại trừ khả năng nhiễm SARS-CoV-2.

\item

  \textbf{Các kết quả ``dương tính'' có thể được Washington State Department of
  Health xác thực}. Quý vị nên thảo luận với nhà cung cấp dịch vụ chăm sóc sức
  khỏe của quý vị về kết quả này nếu quý vị có thắc mắc.

\item{
  \textbf{Nếu xét nghiệm của quý vị cho ra kết quả là ``Không phát hiện'' thì
  điều đó có nghĩa là không phát hiện SARS-CoV-2 trong phòng xét nghiệm của
  chúng tôi.}

  Không có xét nghiệm nào là chính xác tuyệt đối và vẫn có khả năng nhỏ là quý
  vị bị nhiễm vi-rút SARS-CoV-2 gây ra COVID-19 dù kết quả âm tính. Xét nghiệm
  này không sử dụng cách lấy mẫu bệnh phẩm thông thường từ nhà cung cấp dịch vụ
  chăm sóc sức khỏe. Mẫu bệnh phẩm này được gửi qua đường bưu điện hoặc dịch vụ
  giao hàng, chứ không được lấy và vận chuyển trong môi trường y tế. Bất kể kết
  quả xét nghiệm của quý vị ra sao thì tốt nhất quý vị nên ở nhà khi bị bệnh và
  tránh tiếp xúc với người khác. Nếu quý vị gặp các vấn đề sức khỏe khác hoặc bị
  bệnh nặng hay lo lắng về bệnh tình của mình thì quý vị nên liên hệ với nhà
  cung cấp dịch vụ chăm sóc sức khỏe của quý vị. Nếu các triệu chứng vẫn kéo dài
  thì quý vị nên trao đổi với nhà cung cấp dịch vụ chăm sóc sức khỏe về việc
  thực hiện xét nghiệm COVID-19 khác.
}

\item

  \textbf{Nếu xét nghiệm của quý vị cho ra kết quả là ``Chưa có kết luận''} thì
  điều đó có nghĩa là kết quả xét nghiệm SARS-CoV-2 trong phòng xét nghiệm của
  chúng tôi không chắc chắn.

  Không có xét nghiệm nào là chính xác tuyệt đối và vẫn có khả năng rằng kết quả
  chưa có kết luận có thể đồng nghĩa với việc quý vị bị nhiễm vi-rút SARS-CoV-2
  gây ra COVID-19. Xét nghiệm này không sử dụng cách lấy mẫu bệnh phẩm chuẩn từ
  nhà cung cấp dịch vụ chăm sóc sức khỏe. Mẫu bệnh phẩm này được gửi qua đường
  bưu điện hoặc dịch vụ giao hàng, chứ không được lấy và vận chuyển trong môi
  trường chăm sóc sức khỏe. Bất kể kết quả xét nghiệm của quý vị ra sao thì tốt
  nhất quý vị nên ở nhà khi bị bệnh và tránh tiếp xúc với người khác. Nếu quý vị
  gặp các vấn đề sức khỏe khác hoặc bị bệnh nặng hay lo lắng về bệnh tình của
  mình thì quý vị nên liên hệ với nhà cung cấp dịch vụ chăm sóc sức khỏe của quý
  vị. Nếu các triệu chứng vẫn kéo dài thì quý vị nên trao đổi với nhà cung cấp
  dịch vụ chăm sóc sức khỏe về việc thực hiện xét nghiệm COVID-19 khác.

\item

  \textbf{Nếu xét nghiệm của quý vị cho ra kết quả là ``Không xét nghiệm được''}
  thì điều đó có nghĩa hoặc không đủ mẫu bệnh phẩm, hoặc đã xảy ra sai sót trong
  quá trình xét nghiệm.

  Do không thể thực hiện xét nghiệm nên có khả năng là quý vị bị nhiễm vi-rút
  SARS-CoV-2 gây ra COVID-19. Xét nghiệm này không sử dụng cách lấy mẫu bệnh
  phẩm thông thường từ nhà cung cấp dịch vụ chăm sóc sức khỏe. Mẫu bệnh phẩm này
  được gửi qua đường bưu điện hoặc dịch vụ giao hàng, chứ không được lấy và vận
  chuyển trong môi trường y tế. Bất kể kết quả xét nghiệm của quý vị ra sao thì
  tốt nhất quý vị nên ở nhà khi bị bệnh và tránh tiếp xúc với người khác. Nếu
  quý vị gặp các vấn đề sức khỏe khác hoặc bị bệnh nặng hay lo lắng về bệnh tình
  của mình thì quý vị nên liên hệ với nhà cung cấp dịch vụ chăm sóc sức khỏe của
  quý vị. Nếu các triệu chứng vẫn kéo dài thì quý vị nên trao đổi với nhà cung
  cấp dịch vụ chăm sóc sức khỏe về việc thực hiện xét nghiệm COVID-19 khác.

\item

  \textbf{Các xét nghiệm dương tính (cho thấy kết quả là ``PHÁT HIỆN'' hoặc
  ``Phát Hiện Vi-rút'') sẽ được cơ sở xét nghiệm báo cáo lên Sở Y Tế cấp tiểu
  bang và quận có thẩm quyền.} Quý vị KHÔNG cần phải tự báo cáo những kết quả
  này. Sở Y Tế Công Cộng có thể liên hệ với quý vị để thu thập thêm thông tin
  về bệnh tình và những người có thể đã tiếp xúc với quý vị nếu kết quả xét
  nghiệm chủng vi-rút corona mới là dương tính.

\item

  \textbf{Quý vị có thể in ra hoặc gửi kết quả xét nghiệm này cho nhà cung cấp
  dịch vụ chăm sóc sức khỏe của mình qua email.}

\end{itemize}

\bigskip
\textbf{Phương Pháp Xét Nghiệm SARS-CoV-2 (COVID-19) Bằng PCR Định Tính:}

Xét nghiệm này được tạo ra nhằm phát hiện SARS-CoV-2 bằng phương pháp PCR phiên
mã ngược. Xét nghiệm này phụ thuộc vào quá trình khuếch đại các đoạn gen vi-rút
Orf1B và S, trong đó bao gồm việc kiểm soát khuếch đại gen chứng nội RNaseP của
con người. Xét nghiệm này đã được phát triển và Northwest Genomics Center trực
thuộc University of Washington đã xác định được các đặc trưng về tính hiệu quả
của xét nghiệm này. Xét nghiệm này do phòng thí nghiệm phát triển đã được xác
nhận tính hợp lệ và gửi đến Food and Drug Administration của Hoa Kỳ để xin Giấy
Phép Sử Dụng Trong Trường Hợp Khẩn Cấp. Đánh giá độc lập của Cơ Quan Quản Lý
Dược Phẩm (Food and Drug Administration - FDA) về việc xác nhận tính hợp lệ của
xét nghiệm vẫn đang chờ phê duyệt. Phòng thí nghiệm này được chứng nhận là đủ
tiêu chuẩn thực hiện xét nghiệm lâm sàng có độ phức tạp cao trong phòng thí
nghiệm theo Các Quy Định Sửa Đổi Cải Tiến Phòng Thí Nghiệm Lâm Sàng (Clinical
Laboratory Improvement Amendments - CLIA).

Để biết thêm thông tin về nhà cung cấp, vui lòng xem các đường dẫn này do
Washington State Department of Health cung cấp:

\begin{itemize}
\item

  \link{https://www.doh.wa.gov/Portals/1/Documents/1600/coronavirus/COVIDcasepositiveVIETNAMESE.pdf}{Bệnh
  nhân được xác nhận nhiễm hoặc nghi nhiễm COVID-19}

\item

  \link{https://www.doh.wa.gov/Portals/1/Documents/1600/coronavirus/COVIDexposed\_notSickYet\_VI.pdf}{Bệnh
  nhân đã phơi nhiễm với ca được xác nhận nhiễm COVID-19}

\item

  \link{https://www.doh.wa.gov/Portals/1/Documents/1600/coronavirus/COVIDsickNotExposedVietnamese.pdf}{Bệnh
  nhân không phơi nhiễm nhưng có các triệu chứng nhiễm COVID-19}

\end{itemize}

Các nguồn thông tin khác hiện có sẵn trên trang web của Washington State
Department of Health và CDC:

\begin{itemize}
\item

  \url{https://www.doh.wa.gov/Emergencies/NovelCoronavirusOutbreak2020COVID19/vietnamese}

\item

  \url{https://vietnamese.cdc.gov/coronavirus/2019-ncov/index.html}

\end{itemize}

\end{document}
